\section{Introduction}\label{sec:introduction}
Insurance fraud creates financial losses and places demands on claims
investigation. Insurers must identify suspicious claims while allocating
limited review resources \citep{A61,Bolton2002}. This task is difficult
because legitimate claims are heterogeneous and unusual observations are
not necessarily fraudulent.

Machine-learning approaches address this problem through supervised
classification, unsupervised screening, or combinations of the two
\citep{Ngai2011,West2016,Schrijver2024}. Their practical value depends on
both detection performance and the number of claims requiring verification.
Selective investigation also determines which labels become available
for subsequent training.

Sequential Anomaly Detection (SAD) connects initial anomaly screening to
verification and repeated supervised learning. The framework is intended
to support a choice between detection performance and review workload.
We examine this choice through the simple and multilayered SAD structures.

\begin{revised}
The contribution lies in the organization and evaluation of established
components within a verification-driven workflow. Its \emph{sequential}
mechanism incorporates verified records from one batch into training before
the ensemble screens the next batch. Its \emph{multilayer} mechanism
applies an entropy filter to the candidates selected by anomaly-detector
consensus. Updating changes the information available for later prediction;
layering changes which initial cases reach verification. Simple-A,
Simple-B and Multilayer retain the original architectural alternatives:
entropy is absent, contributes a flat vote, or forms a second screening
layer, respectively.

The evaluation compares these three architectures, introduces a frozen-model
control for sequential updating, and applies the same prescribed protocol
to two benchmark datasets. Additional classifiers include an imbalance-aware
forest and a recent fraud-motivated boosting approach. The experiments
explicitly account for review workload and designate fraud as the positive
class. Recorded outcomes simulate expert feedback and are revealed to the
learner only after review selection. SAD uses fixed acquisition rules;
the comparison with reinforcement learning therefore concerns workflow
design rather than a learned reward-maximizing policy.
\end{revised}

Section \ref{SC:LR} positions SAD within the fraud-detection literature.
Section \ref{SC:METHOD} defines its components and update rules;
Section \ref{SC:RESULT} reports the benchmark experiments.
Section \ref{SC:CONCLUSION} discusses the findings and their limitations.
